\section{The original principal cusp lattice and its exact normalization}
\label{sec:principal-cusp-lattices-v047}

\noindent\textbf{Scope.}
This chapter constructs the Hermitian coordinates and the self-dual
lattice that the principal-chart part of ICL retains.  Its generic
endpoint is the full original \(\mathrm{DataX}\) record, on the same
\(\mathrm{Data}\), with the same Hermitian matrix, scalar, arithmetic
group and principal modulus.  On a supplied original uniformization
it gives the exact host denominators and their logarithmic bounds.
Existence of that uniformization and of its principal congruence
inclusion remains an explicit hypothesis.  In particular this chapter
does not assert \(\mathrm{exists\_unf\_hasPrinc}\), construct
\(\mathrm{LawsCore}\), or complete the five remaining ICL spanning
producer groups.

Write \(O=\mathcal O_E\), where \(E\) is a CM number field, and put
\(\bar a=\operatorname{cc}(a)\).  The original convention is
\[
 h_G(x,y)=x^tG\bar y,\qquad
 B_{\xi,G}(x,y)=\operatorname{Tr}_{E/\mathbb Q}
                                      (-\xi h_G(x,y)).       \tag{PCL.1}
\]
Thus \(h_G\) is linear in its first variable.  We write
\(\mathfrak D_E\) for the absolute different and \(aL\) for the
\(O\)-span of products \(\alpha v\), \(\alpha\in a,v\in L\),
where \(a\) is a fractional ideal.  A full lattice is a finitely
generated \(O\)-submodule with the whole \(E\)-span; it need not
have an \(O\)-basis.  The algebraic foundations used here are the
Dedekind ideal laws, the nondegenerate trace pairing for a number
field, and the different as the inverse trace-dual ideal.  Their
specific consequences are derived below.  No local freeness or
principal-ideal hypothesis is imposed on the actual lattices.

% BEGIN PCL1
\paragraph{PCL1. Hermitian coordinates for the original alternating datum.}
\label{claim:PCL1}
Let \(D_w:\mathrm{Datum}(H)\) be the original datum, let
\(r=D_w.p+D_w.q\), and let \(0\ne\xi\in E\) satisfy
\(\bar\xi=-\xi\).  There exist a matrix \(G\in M_r(E)\)
and an \(O\)-linear injection \(\phi:D_w.\mathrm{La}\to E^r\)
such that its image spans \(E^r\), \(\bar G=G^t\),
\(\det G\ne0\), and
\[
 D_w.\mathrm{EF}(x,y)
       =\operatorname{Tr}_{E/\mathbb Q}
                            (\xi h_G(\phi x,\phi y)).       \tag{PCL.2}
\]
This is the literal \(\mathrm{exists\_herm\_coord\_of\_datum}\),
including its sign and its original lattice domain.

\emph{Proof.}
Put \(L=D_w.\mathrm{La}\) and \(b=D_w.\mathrm{EF}\).
The supplied datum says that \(b\) is a perfect alternating
\(\mathbb Z\)-pairing and
\(b(ax,y)=b(x,\bar a y)\) for \(a\in O\).
Perfection makes \(L\) torsion-free over \(\mathbb Z\).
It is also torsion-free over \(O\): a nonzero \(a\in O\)
divides a nonzero rational integer in \(O\), by its integral
equation, so \(ax=0\) would give integer torsion.  Hence
\(L\to V=E\otimes_O L\) is injective.  The original finite
generation and rank fields give \(\dim_E V=r\).

For fixed \(x,y\in L\), the additive functional
\(a\mapsto b(ax,y)\) on \(O\) extends uniquely to a
\(\mathbb Q\)-linear functional on \(E\).  Nondegeneracy
of the trace gives a unique \(t(x,y)\in E\) with
\[
 \operatorname{Tr}(a\,t(x,y))=b(ax,y)\quad(a\in O).
                                                               \tag{PCL.3}
\]
Uniqueness in (PCL.3), additivity, and the adjoint law show
\(t(ax,y)=a\,t(x,y)\),
\(t(x,ay)=\bar a\,t(x,y)\), and additivity in both variables.
Alternation and invariance of the trace under conjugation give
\(t(y,x)=-\overline{t(x,y)}\).  For example, pairing both
sides with every \(a\in O\) reduces this last assertion to
\(b(ay,x)=-b(x,ay)=-b(\bar a x,y)\).
These laws extend \(t\) uniquely to a sesquilinear form on
\(V\), by clearing denominators in each argument.  Then
\(h=\xi^{-1}t\) is Hermitian.  If \(h(v,V)=0\), choose
\(0\ne a\in O\) with \(av\in L\).  Equation (PCL.3) and
perfection give \(av=0\), hence \(v=0\).  Thus \(h\) is
nondegenerate.

Choose an \(E\)-basis of \(V\), take its coordinate map
\(V\simeq E^r\), and let \(G\) be the Gram matrix of \(h\).
Its restriction is \(\phi\); Hermitian symmetry and
nondegeneracy give the matrix assertions.  Taking \(a=1\)
in (PCL.3) proves (PCL.2).  This also covers \(r=0\), with
the empty matrix of determinant one.  For the chart convention
(PCL.1), the input to this theorem is \(-\xi_0\).
\hfill\(\square\)
% END PCL1

% BEGIN PCL2
\paragraph{PCL2. Ideal-valued duals without a principal lattice basis.}
\label{claim:PCL2}
For an \(O\)-lattice \(L\subset E^r\), put
\[
 L^\#=\{x:h_G(x,L)\subset O\}.
\]
For every nonzero fractional ideal \(a\),
\[
 \{x:h_G(x,L)\subset a\}=aL^\#,
 \qquad (aL)^\#=\bar a^{-1}L^\#.                    \tag{PCL.4}
\]
Moreover \(aL\) is a full finitely generated lattice,
\(a(a^{-1}L)=L\), and the product notation is exactly the
span of the products in the original source.

\emph{Proof.}
Since \(aa^{-1}=O\), choose finitely many
\(\alpha_i\in a,\beta_i\in a^{-1}\) with
\(\sum_i\alpha_i\beta_i=1\).  If \(h_G(x,L)\subset a\),
then \(h_G(\beta_i x,L)\subset O\), so
\(x=\sum_i\alpha_i(\beta_i x)\in aL^\#\).
The reverse inclusion follows term by term from linearity.
Next, \(x\in(aL)^\#\) is equivalent to
\(\bar\alpha h_G(x,y)\in O\) for all
\(\alpha\in a,y\in L\), which is equivalent to
\(h_G(x,L)\subset\bar a^{-1}\).  This proves (PCL.4).

Products of finite generating sets generate \(aL\), and any
nonzero \(\alpha\in a\) shows that its \(E\)-span contains
that of \(L\).  Expanding finite sums proves
\((ab)L=a(bL)\).  Together with \(aa^{-1}=O\), this gives
the cancellation identity and the source's displayed span
formula.  All these arguments use invertible ideals, not
generators for them. \hfill\(\square\)
% END PCL2

% BEGIN PCL3
\paragraph{PCL3. The exact different in the perfect trace pairing.}
\label{claim:PCL3}
Assume \(\xi\ne0\) and \(\det G\ne0\).  Define the
trace dual by
\(L^B=\{x:B_{\xi,G}(x,L)\subset\mathbb Z\}\).
Then
\[
 L^B=((\xi)\mathfrak D_E)^{-1}L^\# .              \tag{PCL.5}
\]
If \(L\) is full, the following conditions are equivalent:
the restriction of \(B_{\xi,G}\) is integral and represents
every additive functional \(L\to\mathbb Z\); \(L^B=L\);
and
\[
                         L^\#=(\xi)\mathfrak D_E L. \tag{PCL.6}
\]

\emph{Proof.}
The defining property of the inverse different is
\[
 z\in\mathfrak D_E^{-1}
       \ \Longleftrightarrow\
             \operatorname{Tr}(az)\in\mathbb Z\quad(a\in O).
\]
Because \(\bar a y\in L\) for \(a\in O,y\in L\), the
condition \(x\in L^B\) is equivalent to
\(-\xi h_G(x,y)\in\mathfrak D_E^{-1}\) for every
\(y\in L\).  The sign is a unit and does not change the
ideal.  PCL2 now gives (PCL.5).

The pairing \(B_{\xi,G}\) on the underlying rational vector
space is nondegenerate.  Indeed, if all its values on
\((x,y)\) vanish, also replace \(y\) by \(ay\) for every
\(a\in E\); trace nondegeneracy gives
\(\xi h_G(x,y)=0\), and then \(x=0\).
A full \(O\)-lattice is a full free \(\mathbb Z\)-lattice:
it is finitely generated and torsion-free, and its rational
span is all of \(E^r\).  The latter follows by replacing
an \(O\)-denominator by a nonzero integer multiple of it.
Every additive functional on \(L\) therefore extends to a
rational linear functional, represented uniquely by
\(B_{\xi,G}(x,-)\) for some \(x\in E^r\).  Its integer
values say exactly that \(x\in L^B\).

Consequently \(L^B=L\) gives both integrality and the stated
representation property.  Conversely integrality gives
\(L\subseteq L^B\); for \(x\in L^B\), the representation
property supplies \(x'\in L\) with the same pairing on
\(L\), and nondegeneracy on its rational span gives
\(x=x'\).  Finally cancel the nonzero fractional ideal in
(PCL.5), using PCL2, to obtain (PCL.6).
\hfill\(\square\)
% END PCL3

% BEGIN PCL4
\paragraph{PCL4. The actual fractional-ideal untwisting is self-dual.}
\label{claim:PCL4}
Let \(\Lambda\) carry the perfect integral trace pairing
of PCL3.  Suppose nonzero fractional ideals
\(\mathfrak A,\mathfrak A_b\) satisfy
\[
 \mathfrak A_b=\overline{\mathfrak A},\qquad
             (\xi)\mathfrak A\mathfrak A_b\mathfrak D_E=O.
                                                               \tag{PCL.7}
\]
Then the particular lattice
\[
               L'=\mathfrak A^{-1}\Lambda            \tag{PCL.8}
\]
is full, finitely generated, integral and self-dual for
\(h_G\), and \(\Lambda=\mathfrak A L'\).

\emph{Proof.}
PCL2 gives finite generation, fullness and the last
identity.  Equations (PCL.4), (PCL.6) and (PCL.7) give
\[
 (L')^\#=\overline{\mathfrak A}\,\Lambda^\#
       =\overline{\mathfrak A}(\xi)\mathfrak D_E\Lambda
       =\mathfrak A^{-1}\Lambda=L'.                  \tag{PCL.9}
\]
The inclusion \(L'\subseteq(L')^\#\) is precisely
integrality; its reverse is precisely the source's
\(\mathrm{L'\_dual}\) field.  Thus this is an equality
of the original lattices, with no saturation, denominator
enlargement or change of Hermitian form.
\hfill\(\square\)
% END PCL4

% BEGIN PCL5
\paragraph{PCL5. Construction of the complete original \(\mathrm{DataX}\).}
\label{claim:PCL5}
Let \(X:\mathrm{Data}(E,r)\), and assume exactly the inputs
of \(\mathrm{exists\_dataX\_of\_perfect\_traceForm}\):
(PCL.7) with \(\xi=X.\xi_0\), integrality and perfection
of \(B_{\xi,G}\) on \(X.\Lambda\), a real \(k\le1\)
with \(|\tau(\xi)|=k\) for every complex embedding, and
a positive integer \(g\) whose principal congruence
subgroup in the original unitary stabilizer lies in
\(X.\Gamma\).  There is \(X_x:\mathrm{DataX}(E,r)\) such
that
\[
 X_x.\mathrm{toData}=X,\quad X_x.g=g,\quad
 X_x.L'=\mathfrak A^{-1}X.\Lambda,\quad
 X_x.\Lambda=\mathfrak A X_x.L'.                     \tag{PCL.10}
\]

\emph{Proof.}
The ideal equation forces \(\mathfrak A\) and
\(\mathfrak A_b\) to be nonzero; their given
conjugate-membership condition says exactly that
\(\mathfrak A_b=\bar{\mathfrak A}\).
Every nonzero fractional ideal \(a\) contains a
positive rational integer.  Choose \(0\ne u\in a\),
multiply it by a nonzero integer so that \(u\in O\), and
use the finite additive group \(O/(u)\): its positive
order \(N((u))\) kills the class of \(1\), so
\(N((u))\in(u)\subset a\).  Choose in this way
\(N_A>0\) in \(\mathfrak A^{-1}\) and \(N_C>0\) in
\(\mathfrak A\).  Then
\[
                N_A\Lambda\subset L',\qquad
                            N_C L'\subset\Lambda.    \tag{PCL.11}
\]
There is a complex embedding of \(E\); its injectivity
and \(\xi\ne0\) give \(k>0\).  Set \(k_0=2k\).
It satisfies \(0<k_0\le2\) and the exact field
\(|\tau(\xi_0)|=k_0/2\) on the source CM type.

Keep the entire record \(X\).  Supply \(\mathrm{EF\_int}\)
and \(\mathrm{EF\_perf}\) by the two given trace-form
properties.  Supply \(L'\), its finite generation,
\(\mathrm{L'\_int}\), and \(\mathrm{L'\_dual}\) by PCL4.
Supply the two positive integers and both comparison
inclusions by (PCL.11).  Set its modulus to the given
\(g\); positivity and the \(\mathrm{cong}\) field are
exactly the assumed principal congruence inclusion.
These are every additional field of \(\mathrm{DataX}\).
The last equality in (PCL.10), expanded using PCL2, is
the literal span of \(\alpha u\) in the original endpoint.
In particular the construction changes no scalar,
matrix, character, group, domain or pre-existing metric.
\hfill\(\square\)
% END PCL5

% BEGIN PCL6
\paragraph{PCL6. Perfection on every component of the supplied uniformization.}
\label{claim:PCL6}
For any original \(U:\mathrm{Unf}(X_o,c,\sigma,r)\)
satisfying \(\mathrm{IsUnf}(U)\), each component
\(X=U.X(k)\) has the integral perfect trace pairing
in PCL5.  This proves the full
\(\mathrm{traceForm\_integral\_perfect\_of\_isUnf}\).

\emph{Proof.}
Unpack the actual \(\mathrm{IsUnf}\) witness.  It supplies
a family of original alternating data \(D_i\), a component
index \(i=i(k)\), and an \(O\)-linear equivalence
\(\phi_k:D_i.\mathrm{La}\simeq\Lambda_k\) such that
\[
        D_i.\mathrm{EF}(a,b)
             =B_{\xi_k,G_k}(\phi_k a,\phi_k b).       \tag{PCL.12}
\]
For two vectors in \(\Lambda_k\), apply \(\phi_k^{-1}\)
and use the integer-valued left side.  For an additive
\(\psi:\Lambda_k\to\mathbb Z\), perfection of
\(D_i.\mathrm{EF}\) represents \(\psi\circ\phi_k\)
by some \(a\in D_i.\mathrm{La}\).  Equation (PCL.12)
then represents \(\psi\) by \(\phi_k a\in\Lambda_k\).
This argument uses the lattice equivalence appearing in
the supplied uniformization; no independently chosen
lattice or rational isometry is substituted.
\hfill\(\square\)
% END PCL6

% BEGIN PCL7
\paragraph{PCL7. The two actual quadratic ideal identities.}
\label{claim:PCL7}
For the original host, let \(K=\mathbb Q(s)\),
\(s=\sqrt{-d}\), let \(C=\mathrm{ccK}\ne0\), and
let \(C_E=CO_E\).  The original prime and congruence
conditions give \(d>0\) squarefree and \(d\equiv3\pmod4\).
Then
\[
       \mathfrak D_K O_E=(s),\qquad
          C_E\overline{C_E}=(N_K C).                \tag{PCL.13}
\]
These are the literal \(\mathrm{DiffKPrincipal}\) and
\(\mathrm{CcNorm}\) suppliers on the original host.

\emph{Proof.}
The host conjugation is the CM conjugation used in
(PCL.1): the host's embedding law and the CM embedding
law both send it under \(\phi_0\) to complex conjugation,
so injectivity identifies them on all of \(E\).
The integer ring of \(K\) is
\(\mathbb Z[\omega]\), \(\omega=(1+s)/2\).
Here is the elementary integral-coordinate check.  If
\(x=a+bs\) is integral, its trace gives \(2a\in\mathbb Z\)
and its norm gives \(a^2+db^2\in\mathbb Z\).
For a rational denominator of \(b\) in lowest terms,
these conditions imply its square divides \(4d\);
squarefreeness of \(d\) makes that denominator divide
\(2\).  Writing \(x=(A+Bs)/2\), the norm condition is
\(A^2+dB^2\equiv0\pmod4\), so \(A\equiv B\pmod2\).
Conversely \(\omega\) satisfies the monic equation
\(T^2-T+(d+1)/4=0\).  This proves the stated integer ring.

Its trace dual is \(s^{-1}O_K\).  Indeed the trace matrix
of \(1,\omega\) has determinant \(-d\), so the trace dual
has index \(d\) over \(O_K\).  Directly,
\(\operatorname{Tr}(1/s)=0\),
\(\operatorname{Tr}(\omega/s)=1\), and
\(\operatorname{Tr}(\omega^2/s)=1\); hence
\(s^{-1}O_K\) lies in the trace dual.  It also has index
\(|N_{K/\mathbb Q}(s)|=d\), proving equality.
Inverting gives \(\mathfrak D_K=(s)\).

For the norm identity, factor \(C\) into prime ideals.
Over a split rational prime the two primes are exchanged
by conjugation and their product is \((p)\); each has
norm \(p\).  At an inert prime the single prime is
\((p)\), with norm \(p^2\), so its product with its
conjugate is \((p^2)\).  At a ramified prime its square
is \((p)\) and its norm is \(p\).  These three cases
give \(P\bar P=(NP)\) for every prime factor.  Products
give \(C\bar C=(N_K C)\).  Extending both identities
to \(E\) gives (PCL.13); host conjugation restricts to
the nontrivial conjugation of \(K\), since it sends
\(s\) to \(-s\). \hfill\(\square\)
% END PCL7

% BEGIN PCL8
\paragraph{PCL8. The original principal ideal and scalar normalization.}
\label{claim:PCL8}
Write \(\mathfrak b\) for the stored half-different, so
\(\mathfrak D_{E/K}=\mathfrak b\bar{\mathfrak b}\), and put
\[
 N=N_K C,\qquad
 \xi_H=-\frac1{Ns},\qquad
 \mathfrak A=C_E\mathfrak b^{-1},\qquad
 k_H=\frac1{N\sqrt d}.                              \tag{PCL.14}
\]
These are exactly the source's \(\xi_0,\mathrm{AA}\) and
\(\xi R\), with \(\mathrm{AAbar}=\bar{\mathfrak A}\).
They satisfy
\[
 (\xi_H)\mathfrak A\bar{\mathfrak A}\mathfrak D_E=O,
 \qquad \bar\xi_H=-\xi_H,
 \qquad |\tau(\xi_H)|=k_H\quad\hbox{for every }\tau,
 \qquad 0<k_H\le1.                                 \tag{PCL.15}
\]

\emph{Proof.}
The different in a tower gives
\(\mathfrak D_E=\mathfrak b\bar{\mathfrak b}(s)\).
Both half-different factors are nonzero because the
different is nonzero.  Apply PCL7 and cancel these
invertible ideals:
\[
 (\xi_H)(C_E\mathfrak b^{-1})
       (\bar C_E\bar{\mathfrak b}^{-1})
       \mathfrak b\bar{\mathfrak b}(s)
       =(\xi_H Ns)=(-1)=O.                          \tag{PCL.16}
\]
This is a proof of the actual \(\mathrm{FactPrincipalDatum}\),
not an additional principal-ideal assumption on
\(\mathfrak A\).  Conjugation of the definition gives
\(\bar\xi_H=-\xi_H\).  Every complex embedding sends
\(s\) to a root of \(z^2=-d\), whose absolute value is
\(\sqrt d\).  Hence its absolute value on \(\xi_H\)
is the same positive \(k_H\).  Since \(N\ge1,d\ge1\),
one has \(k_H\le1\).  At the source's fixed embedding
this equality is the definition of \(\xi R\).
\hfill\(\square\)
% END PCL8

% BEGIN PCL9
\paragraph{PCL9. The exact two host denominators.}
\label{claim:PCL9}
For any lattice \(\Lambda=\mathfrak A L'\) with the
original \(\mathfrak A\) of (PCL.14), set
\[
             N_A=N_E\mathfrak b,
             \qquad N_C=N_K C.                     \tag{PCL.17}
\]
Both are positive and
\[
               N_A\Lambda\subset L',\qquad
                           N_C L'\subset\Lambda.    \tag{PCL.18}
\]
These integers work for all components and embeddings;
they are the literal \(\mathrm{NA}\) and \(\mathrm{NC}\)
in the quantitative original endpoint.

\emph{Proof.}
For a nonzero integral ideal \(I\), the additive group
\(O/I\) has \(NI\) elements; thus \(NI\in I\).
In particular \(N_A\in\mathfrak b\), so
\(N_A\mathfrak A\subset C_E\subset O\).  Multiplying
\(\Lambda=\mathfrak A L'\) proves the first inclusion.
Also \(N_C\in C\), hence \(N_C\in C_E\).
Since \(\mathfrak b\subset O\),
\(N_C\mathfrak b\subset C_E\); equivalently
\(N_C\in C_E\mathfrak b^{-1}=\mathfrak A\).
This proves the second inclusion.  Nonzero ideals have
nonzero finite norm, so both integers are positive.
No ideal is principalized, and no integer depending on
the component has replaced either value.
\hfill\(\square\)
% END PCL9

% BEGIN PCL10
\paragraph{PCL10. The full half-different logarithmic bound.}
\label{claim:PCL10}
For the original host,
\[
 0<N_E\mathfrak b,\qquad
 (N_E\mathfrak b)^2\le|\operatorname{disc}(E)|,
 \qquad
 2\log N_E\mathfrak b\le\log|\operatorname{disc}(E)|.
                                                               \tag{PCL.19}
\]
This supplies \(\mathrm{HalfDiffNormSt}\), including its
positivity.  With the exact source definitions
\(f_E=[E:\mathbb Q]/2\) and
\(\operatorname{rd}(E)=|\operatorname{disc}(E)|^{1/[E:\mathbb Q]}\),
it gives
\[
        \log N_A\le f_E\log\operatorname{rd}(E).     \tag{PCL.20}
\]

\emph{Proof.}
Conjugation gives a ring isomorphism
\(O_E/\mathfrak b\simeq O_E/\bar{\mathfrak b}\), so their
finite cardinalities agree.  Multiplicativity of ideal
norm therefore gives
\(N_E\mathfrak D_{E/K}=(N_E\mathfrak b)^2\).
The different-in-a-tower identity writes
\(\mathfrak D_E\) as \(\mathfrak D_{E/K}\) times a
nonzero integral ideal.  The norm of that latter ideal
is a positive integer, and
\(N_E\mathfrak D_E=|\operatorname{disc}(E)|\).  This
proves the first two assertions.  Take the increasing
logarithm of the positive quantities to obtain the third.

A CM field has its complex embeddings paired by
fixed-point-free conjugation; thus its positive degree
is even and the source's natural-valued division by two
is exactly \([E:\mathbb Q]/2\).  Taking the logarithm
of the positive real power in the definition of
\(\operatorname{rd}\) shows
\(f_E\log\operatorname{rd}(E)=\tfrac12\log|\operatorname{disc}(E)|\).
This proves (PCL.20), with no hidden factor of two.
It also implies the weaker original bound
\(N_A\le|\operatorname{disc}(E)|\), since the latter
is a positive integer.
\hfill\(\square\)
% END PCL10

% BEGIN PCL11
\paragraph{PCL11. The literal normalized \(\mathrm{DataX}\) on every supplied chart.}
\label{claim:PCL11}
Fix the original \(D\) and \(T:\mathrm{Core}(D)\),
\(h:\mathrm{LawsCore}(T)\), a row \(j\) with
\(\mathrm{InR}(D,j)\), and any good embedding \(\sigma\).
Assume the original
\(\mathrm{HasPrinc}(T.U(j,\sigma),T.Q.l.M)\).
For every component \(k\), there is \(X_x\) with
\[
 \begin{gathered}
 X_x.\mathrm{toData}=(T.U(j,\sigma)).X(k),\qquad
 X_x.g=T.Q.l.M,\qquad X_x.k_0=2\xi R(D),\\
 X_x.L'=\mathrm{AA}^{-1}X_x.\Lambda,\qquad
 X_x.\Lambda=\mathrm{AA}\,X_x.L',\\
 X_x.N_A=N_E\mathfrak b,\qquad
 X_x.N_C=N_K\mathrm{ccK},\qquad
 X_x.N_A\le|\operatorname{disc}(E)|.
 \end{gathered}                                                   \tag{PCL.21}
\]
The fractional-ideal equality is the original span of
\(\alpha u\).  In particular this proves the entire
\(\mathrm{exists\_dataX\_of\_hasPrinc}\), with equality
in its advertised \(N_C\) bound.

\emph{Proof.}
The actual \(h.\mathrm{unf}\) field supplies PCL6, so
the trace pairing on this precise component is integral
and perfect.  Write its scalar as \(\xi\).
The field \(h.\mathrm{xiAbs}\) says
\(|\tau(\xi)|=\xi R(D)=|\tau(\xi_H)|\) at every
embedding, while \(\bar\xi=-\xi\) is part of its
\(\mathrm{Data}\).  At the one fixed embedding \(\phi_0\),
both values are purely imaginary and have equal absolute
value; hence one is the other or its negative.
Injectivity of \(\phi_0\) gives \(\xi=\xi_H\) or
\(\xi=-\xi_H\).  Their principal ideals coincide,
so PCL8 gives (PCL.7) for this actual component.

Apply PCL5 with \(\mathfrak A=\mathrm{AA}\),
\(k=\xi R(D)\), and \(g=T.Q.l.M\).
Its positive-modulus assumption follows from the
original level field \(M\ge3\); its congruence assumption
is precisely the given \(\mathrm{HasPrinc}\), with the
matrix acting on vectors on the left as in (PCL.1).
Use the prescribed integers in PCL9 instead of the
arbitrary integers in the proof of PCL5.  All record
fields still hold, now with (PCL.21); PCL10 gives the
claimed norm bound.  No change to \(T.U\), its component
index, \(G\), \(\xi\), \(\Gamma\), character or modulus
has occurred.  The argument works for every row and
every good embedding quantified in \(\mathrm{LawsCore}\),
not merely its distinguished complex reading.
\hfill\(\square\)
% END PCL11

% BEGIN PCL12
\paragraph{PCL12. The exact all-row model-lattice endpoints and remaining boundary.}
\label{claim:PCL12}
For the original core, retain \(\mathrm{GamAt}(T)\) and
\(\mathrm{LawsCore}(T)\).  Then the full
\(\mathrm{dataXOfModel\_of\_gamAt}\) holds: for every
row \(j\) satisfying \(\mathrm{InR}(D,j)\), every good
embedding and every component, there is
the unchanged \(\mathrm{Data}\) extended to \(\mathrm{DataX}\),
with \(g=T.Q.l.M\) and the same \(\mathrm{AA}\)-span equality.
For the original top record, the complete
\(\mathrm{DataXOfModelSt}(T,j)\) also holds.  In particular,
for every original \(\mathrm{large}_2\) and row weight \(w\),
under its \(\mathrm{SupBoundPageRows}\) and
\(2\le\mathrm{large}_2.\mathrm{sc}.L\) hypotheses, the
same chosen \(X_x\) satisfies
\[
 \log X_x.N_A\le f_E\log\operatorname{rd}(E),\qquad
 \log X_x.N_C\le2\,\mathrm{large}_2.\mathrm{sc}.L.     \tag{PCL.22}
\]

\emph{Proof.}
The source has a core version and a top version of
\(\mathrm{GamAt}\).  Each says, literally, for every
row, good embedding, component and unitary stabilizer
element \(\gamma\),
\[
       (\gamma-1)\Lambda\subset (T.Q.l.M)\Lambda
                      \ \Longrightarrow\ \gamma\in\Gamma.
                                                               \tag{PCL.23}
\]
Thus for the specified row and embedding it supplies
the exact \(\mathrm{HasPrinc}\) required by PCL11.
That claim gives the first endpoint, with all its
quantifiers.  Passing from \(\mathrm{Top}\) to its
\(\mathrm{Core}\) changes none of these data.

Choose the lattice by (PCL.8) and the integers by
(PCL.17), before any weight or \(\mathrm{large}_2\).
PCL10 gives the first inequality in (PCL.22).
The fourth original conjunct of
\(\mathrm{SupBoundPageRows}\) is exactly
\(\log N_K\mathrm{ccK}\le2\,\mathrm{large}_2.\mathrm{sc}.L\),
which gives the second.  All other conclusions of
\(\mathrm{DataXOfModelSt}\) are already (PCL.21).
Its extra size and weight hypotheses may therefore
be retained verbatim without changing the construction.
This also shows that the conclusion is uniform in
every positive original weight, and in both original
level choices whenever those supplied cores are used.
\hfill\(\square\)

The new result is the actual lattice and normalization
producer on a supplied principal chart.  The actual
existence of a common \(\mathrm{IsUnf}\) family with
\(\mathrm{HasPrinc}\), its framed-point normalization,
and its fixed metric dictionary remains separate.
So do the finite \(\mathrm{IsSystem}\), Gaussian
reproduction, short-rule gluing and ramified-jet inputs
retained by ICL.  This chapter does not infer full
\(\mathrm{LatFull}\), a higher \(\mathrm{CoversAt}\),
all-tuple class-sum or phase identities, or any compound
common-family closure from the constructed \(\mathrm{DataX}\).
% END PCL12
